\begin{helper}
Let \(C\) be a nonempty finite set of candidates, let \(p\in[0,1]\), and let
\(\lambda\) be a measure on the space of numerical candidate-priority vectors
\(\pi:C\to\mathbb R\).  Suppose \(\lambda\)-almost every priority vector lies
in the unit cube, that is \(0\le \pi(x)\le 1\) for every \(x\in C\).
For every priority vector \(\pi\), let \(\eta_\pi\) be a one-common-blocker
law and let \(\theta_\pi\) be a private-blocker product law indexed by \(C\).
Assume that, for every \(\pi\), the laws \(\eta_\pi\) and \(\theta_\pi\)
satisfy the total-mass, active-unit, lower-rectangle, and product-cylinder
hypotheses appearing in
\(\noderef{SamplingFiniteProductCandidateReplacementStep}\), with the same
activation parameter \(p\).  Then the integrated probability of candidate
survival after replacing the common blocker by private blockers is monotone:
\[
\int
\eta_\pi\{\neg(a=1,\ q\in[0,1],\ \pi(x)<q\hbox{ for every }x\in C)\}
\,d\lambda(\pi)
\]
is at most
\[
\int
\theta_\pi\{\exists x\in C:
\neg(a_x=1,\ q_x\in[0,1],\ \pi(x)<q_x)\}
\,d\lambda(\pi).
\]
\end{helper}
\begin{proof}
For \(\lambda\)-almost every vector \(\pi\), the support hypothesis gives
\(0\le \pi(x)\le1\) for every \(x\in C\).  Fix such a vector \(\pi\).  The
activation parameter \(p\) satisfies \(0\le p\le1\) by hypothesis, and the
families \(\eta_\pi\) and \(\theta_\pi\) satisfy exactly the one-common-blocker
and private-blocker product-law hypotheses required by
\(\noderef{SamplingFiniteProductCandidateReplacementStep}\).  Applying that
node at this numerical vector \(\pi\) gives the pointwise inequality
\[
\eta_\pi\{\neg(a=1,\ q\in[0,1],\ \pi(x)<q\hbox{ for every }x\in C)\}
\le
\theta_\pi\{\exists x\in C:
\neg(a_x=1,\ q_x\in[0,1],\ \pi(x)<q_x)\}.
\tag{1}
\]
This use is pointwise in the actual numerical priority vector; no event is
conditioned on the value of \(\pi\), and no finite order chamber is treated as
though it determined the numbers \(\pi(x)\).

The nonnegative Lebesgue integral is monotone under almost-everywhere
inequality of its integrands.  Since (1) holds for \(\lambda\)-almost every
\(\pi\), integrating both sides with respect to \(\lambda\) gives precisely
the displayed integrated inequality.  Thus the one-step blocker comparison
continues to hold after averaging over the candidate-priority vector itself.
\end{proof}
